\chapter{Prerequisites for Objective}

\section{Physically Unclonable Function}

A Physically Unclonable Function (PUF) is a one-way physical function that
generates a unique response $R$ for a given challenge $C$ from the challenge
space $\mathcal{C}$.
PUFs exploit manufacturing variations in integrated circuits that are practically
impossible to replicate.
Various types of PUFs have been proposed, including SRAM PUFs, Arbiter PUFs, and
Ring Oscillator (RO) PUFs.
In this work, a Ring Oscillator PUF is used due to its higher execution speed and
lower overhead.

A function is considered a valid and secure PUF if it satisfies the following
properties:

\begin{itemize}
  \item \textbf{Evaluatable:} A legitimate entity can easily compute a response
    $R \in \mathcal{R}$ for any given challenge $C \in \mathcal{C}$.

  \item \textbf{Unpredictability:} An adversary cannot predict the response $R$ for
    any challenge $C$ without physical access to the PUF.

  \item \textbf{Uniqueness:} For two different PUF instances, $\mathsf{PUF}_1$ and
    $\mathsf{PUF}_2$, the probability that
    $\mathsf{PUF}_1(C) = \mathsf{PUF}_2(C)$ for the same challenge $C$ is
    negligible.

  \item \textbf{Reliability:} For a fixed challenge $C$, the PUF response
    $\mathsf{PUF}(C) = R$ remains stable over time and environmental variations.

  \item \textbf{One-wayness:} Given a response $R$, it is computationally infeasible
    to determine the corresponding challenge $C$ or invert the PUF.
\end{itemize}

These properties make PUFs suitable as hardware-bound secrets for key generation,
device binding, and authentication in resource-constrained IoT devices.

\section{Hash-based One-Way Accumulator}

A one-way hash accumulator, proposed as an efficient alternative to digital
signatures in authentication protocols, provides a compact, non-invertible
representation of a set of elements.
In such an accumulator, items are combined using a quasi-commutative operator,
so the final accumulated value does not depend on the order of insertion.

Formally, a function
\[
  f : A \times B \to A
\]
is said to be quasi-commutative if, for all $a \in A$ and for all $b, c \in B$,
\begin{equation}
  f(f(a,b),c) = f(f(a,c),b).
\end{equation}

Nyberg introduced a fast accumulated hashing method suitable for lightweight systems.
Let $N$ denote the maximum number of accumulable items, and let
$h : \{0,1\}^* \to \{0,1\}^\ell$ be a cryptographically secure hash function.
Given a set of items $Y = \{y_1, y_2, \ldots, y_m\}$, $m \le N$, each item is
first hashed as $Y_i = h(y_i)$, $i = 1, 2, \ldots, m$.
The accumulated hash is then computed as:
\begin{equation}
  T_{\mathsf{Acc}} = H(s, Y) = s \bigodot_{i=1}^{m} h(y_i),
\end{equation}
where $s$ is an initialisation seed and $\bigodot$ represents the bitwise AND
operator.
Since the AND operator is commutative and the hash function is one-way, the
accumulator inherits quasi-commutative and one-way properties.

\textbf{Membership Verification:} To verify whether an item $y_i$ belongs to the
accumulated set, compute $Y_i = h(y_i)$, and check whether
$Y_i \bigodot T_{\mathsf{Acc}} = T_{\mathsf{Acc}}$.

The security of the accumulator depends on the difficulty of forging a value $Y_i$
that passes verification.
For a security parameter $t$, the required length $r$ of the accumulated hash is
given by $r = N \times e \times t$, where $e$ is Euler's (Napier's) constant.

\section{Core Concept of Anonymity Preserving}

In IoT security, anonymity is defined by two key properties:
\textbf{Untraceability} (concealing the sender's real identity) and
\textbf{Unlinkability} (preventing the correlation of multiple sessions to a single
device).
While simply encrypting a static ID protects the credential itself, it does not
prevent traffic analysis; if the encrypted value remains constant, an adversary can
track the device's activity patterns.

To achieve robust anonymity, modern schemes utilise \textbf{Dynamic Pseudonyms}.
Instead of a static identifier, the device generates a session-specific alias
($\mathit{PID}$) using a shared secret or random nonce, e.g.,
$\mathit{PID}_{i+1} = H(\mathit{ID} \| \mathit{Nonce}_i)$.
To an external observer, each authentication request appears as a fresh, random
string of bits, completely uncorrelated with previous sessions.

A critical challenge when using dynamic pseudonyms is \textbf{desynchronisation}.
If the message carrying the new nonce is lost in transit, the device retains the old
nonce while the server has already advanced to the new one.
Their pseudonym computations then diverge, and all subsequent authentications fail.
To handle this gracefully, a dual-state storage approach retains both the current and
the immediately previous state at the server, allowing recovery without re-enrolment.
